
But every polynomial @@M0@@ is continuous for the topology induced by @@M1@@. Thus for sufficiently large @@M2@@ one has
@@M3@@
By definition, however, @@M4@@ is an integer for every @@M5@@. It follows that @@M6@@ for sufficiently large @@M7@@. Thus the polynomial @@M8@@ is constant, and in particular one obtains @@M9@@, i.e. @@M10@@.

The two cases a) and b) above show that the restriction of @@M11@@ to @@M12@@ is constant. We shall now show that @@M13@@ is constant on each coset in @@M14@@ modulo @@M15@@. For every @@M16@@, denote by @@M17@@ the map defined by
@@M18@@
The map @@M19@@ is a polynomial map from @@M20@@ to @@M21@@ that factors through @@M22@@, hence is constant by the preceding argument, whence the assertion that @@M23@@ is constant on the cosets in @@M24@@ modulo @@M25@@. It follows that one can write @@M26@@, where @@M27@@ is a polynomial map that factors through @@M28@@. But by hypothesis @@M29@@ is a torsion group, so @@M30@@ is constant. Thus @@M31@@ is constant, Q.E.D.

\sgaStableTarget{sga6.e10.lem.2-1-4}\textbf{Lemma 2.1.4.} \uline{Let} @@M32@@ \uline{be a unital commutative ring and} @@M33@@ \uline{a nilpotent ideal of} @@M34@@\uline{. Let} @@M35@@ \uline{be the multiplicative subgroup of} @@M36@@ \uline{consisting of the invertible elements of} @@M37@@ \uline{of the form} @@M38@@ \uline{with} @@M39@@\uline{, and let} @@M40@@ \uline{be the natural injection. Then the composite map}

